\documentclass[10pt,a4paper]{amsart}
\usepackage[margin=19mm]{geometry}
\usepackage{amsmath,amssymb,mathtools}
\usepackage[scr=rsfs,cal=euler]{mathalfa}
\usepackage{xfrac}
\usepackage[unicode,hidelinks,bookmarks=false]{hyperref}
\newcommand{\ie}{i.e.\ }
\newcommand{\cf}{cf.\ }
\newcommand{\resp}{respectively\ }
\newcommand{\Subsection}{\subsection}
\providecommand{\nobreakdash}{\nobreak\hbox{-}}
\setlength{\emergencystretch}{2em}
\multlinegap0pt
\usepackage{bm}
\usepackage{bbm}
\usepackage{dsfont}
\usepackage{xspace}
\usepackage{cancel}
\usepackage{mathtools}
\usepackage{cleveref}
\usepackage{amsthm}
\makeatletter
\@ifundefined{theorem}{\newtheorem{theorem}{Theorem}}{}
\@ifundefined{corollary}{\newtheorem{corollary}[theorem]{Corollary}}{}
\@ifundefined{lemma}{\newtheorem{lemma}[theorem]{Lemma}}{}
\makeatother
\newcommand{\X}{\mathcal X}
\title[Pinsker's inequality for adapted total variation]{Pinsker's inequality for adapted total variation}
\author{Mathias Beiglb\"ock, Markus Zona}
\date{Source: arXiv:2506.22106v1 (CC BY 4.0)}
\begin{document}
\maketitle

\noindent\emph{Source and licence.} Pinsker's inequality for adapted total variation. Version arXiv:2506.22106v1 (CC BY 4.0), distributed under the Creative Commons Attribution 4.0 International licence, as recorded in the arXiv licence metadata for this exact version and shown on the abstract page https://arxiv.org/abs/2506.22106v1. Full rights notice: licenses/arxiv-2506.22106-CC-BY-4.0.txt.\par\smallskip

\noindent\textit{Editorial scope.} Counted unit: the Corollary whose source label is \texttt{cor:PinskerTight} in the article (source0.tex lines 386--388, source label \texttt{cor:PinskerTight}), delivered together with the article's own complete original proof of it (source0.tex lines 390--410, one \texttt{proof} environment of 21 lines). No mathematics is written, repaired or re-derived: statement and proof are the article's own text line for line. Required context retained uncounted: thm:MainTheorem (lines 247--251); le:ATVsimple (lines 287--297); tightpinsker (lines 373--375). Every definition, display and label of those lines is retained unchanged. Omitted from this package and not redistributed: 1--246, 252--286, 298--372, 376--385, 389--389, 411--423. Nothing inside a retained excerpt is omitted or altered: every retained body is delivered exactly as the article's lines are written, so no omission receipt of any kind accompanies this package. Citations: the retained excerpts carry no citation command at all, so no bibliography entry of the article is transcribed and no reference is redistributed. Reference adaptation: none was needed and none was made; every reference inside the delivered unit resolves inside it, so no unresolved reference or citation is produced. Printed slips kept literal: no printed word, symbol, hypothesis or number of the counted statement or of its proof was altered. Layout and class: the article's own class is \texttt{amsart} with packages including hyperref, cleveref, amssymb, amsmath, mathtools, enumerate, color, mathrsfs, enumitem, graphicx, float, tikz, setspace, geometry; the wrapper is the lane's pinned \texttt{amsart} editorial wrapper at 10pt on A4, which restates the theorem-like environments the retained text needs and adds the article's own macro definitions that the retained text needs (\texttt{\textbackslash X}), with extra wrapper packages (bm, bbm, dsfont, xspace, cancel, mathtools, cleveref). The delivered numbering is the unit's own. Assets: no retained span contains a figure environment, a graphics-inclusion command, a PDF-page inclusion, a table or a TikZ picture; no image file is copied into this package, the package declares no asset dependency and no image file is redistributed with it. Licence: version arXiv:2506.22106v1 of this article is distributed under the Creative Commons Attribution 4.0 International licence, as recorded in the arXiv licence metadata for this exact version and shown on the abstract page https://arxiv.org/abs/2506.22106v1. A full rights notice is delivered at licenses/arxiv-2506.22106-CC-BY-4.0.txt. Characters: every retained excerpt is pure ASCII, so the delivered engine, XeLaTeX with its default Latin Modern text font, reports no missing character.
\begin{theorem}\label{thm:MainTheorem}
  Assume that $\mu, \nu$ are probabilities on the Polish space $\X_1\times \ldots \times \X_n$. Then
  \begin{align}\label{eq:AP} 
  ATV(\mu, \nu)\leq \sqrt{n}\sqrt{2H(\mu|\nu)}.\end{align}
\end{theorem}

\begin{lemma}\label{le:ATVsimple}
    %Suppose $X_1,\ldots, X_n$ are polish spaces and let $\mu, \nu \in \mathcal{P}(X)$, where $X:= X_1 \times \ldots \times X_n$. Then we have the following representation of the adapted total variation between measures $\mu,\nu \in \mathcal{P}(X)$
   The adapted total variation distance can be written as
    \begin{align*}
        ATV(\mu,\nu) & = TV(\mu_1,\nu_1) + \int_{\X_1} TV(\mu^{x_1}, \nu^{x_1}) \, \mathrm{d}(\mu_1 \wedge \nu_1)(x_1) + \ldots + \\
        & \int_{\X_{n-1}} TV(\mu^{x_{1:(n-1)}},\nu^{x_{1:(n-1)}}) \, \mathrm{d}(\mu^{x_{1:(n-2)}} \wedge \nu^{x_{1:(n-2)}})(x_{(n-1)}) \ldots \mathrm{d}(\mu_1 \wedge \nu_1)(x_1).
    \end{align*}
    Moreover the infimum in the definition of $ATV$ is attained.
    Indeed, any coupling $\pi$ for which $\pi_1$ and all $ \pi^{x_{1:k},y_{1:k}}, k< n$ put the maximal possible mass on the diagonal is optimal. 
%    holds.
\end{lemma}

\begin{lemma}\label{tightpinsker}
    The Pinsker inequality is tight.
\end{lemma}

\begin{corollary}\label{cor:PinskerTight}
    The adapted Pinsker inequality $ATV(\mu,\nu) \leq \sqrt{n} \sqrt{2 H(\mu|\nu)}$ is tight for every $n \in \mathbb{N}$. % Ich nutze \mathbb{N}={1,2,3,...}
\end{corollary}

\begin{proof}
    The case $n=1$ is exactly \Cref{tightpinsker}, thus we assume that $n \ge 2$ for the rest of the proof. Consider the measurable space $(\Omega:=\{0,1\}^n,2^\Omega)$ and measures $\mu_i=(\frac{1}{2}+\varepsilon)\ \delta_0 + (\frac{1}{2}-\varepsilon)\ \delta_1$ and $\nu_i = \frac{1}{2}\delta_0+\frac{1}{2}\delta_1$ on $\{0,1\}$ for every $i\in \{1,\ldots,n\}$ as in \Cref{tightpinsker}. Further, define the product measures $\mu := \bigotimes_{i=1}^n \mu_i$ and $\nu := \bigotimes_{i=1}^n \nu_i$. Note that the disintegration of $\mu$ is given by $\mu=\underbrace{\mu_1 \otimes \cdots \otimes \mu_1}_\text{$n$-times}$ 
    and analogously for $\nu$.

    In order to determine $ATV(\mu,\nu)$ we use the decomposition from \Cref{le:ATVsimple} and the fact, that $|\mu_1 \wedge \nu_1| = 1-\varepsilon$ to obtain
    \begin{align*}
        & \phantom{=} ATV(\mu,\nu) \\
        & = TV(\mu_1,\nu_1)+\int_{\X_1} TV(\mu_1,\nu_1)+\ldots +\int_{\X_{n-1}} TV(\mu_1,\nu_1) \ \mathrm{d}(\mu_1 \wedge \nu_1)(x_{n-1}) \ldots \mathrm{d}(\mu_1 \wedge \nu_1)(x_1) \\
        & = 2\varepsilon \sum_{k=0}^{n-1}(1-\varepsilon)^k  = 2\varepsilon \frac{1-(1-\varepsilon)^n}{\varepsilon}  = 2-2(1-\varepsilon)^n.
    \end{align*}
    Since $\mu$ and $\nu$ are product measures, it follows similarly as in the proof of \Cref{tightpinsker} that $H(\mu,|\nu) = n H(\mu_1|\nu_1) = n (2\varepsilon^2+o(\varepsilon^2))$. Due to \Cref{thm:MainTheorem} we  already know that $\frac{ATV(\mu,\nu)^2}{nH(\mu|\nu)}\leq2$ or equivalently $\frac{(2-2(1-\varepsilon)^n)^2}{n^2(2\varepsilon^2+o(\varepsilon^2))} \leq 2$. We conclude the proof by showing that the left hand side of the inequality converges to $2$ as $\varepsilon$ approaches $0$:
    \begin{align*}
        & \phantom{=} \lim_{\varepsilon \searrow 0} \frac{(2-2(1-\varepsilon)^n)^2}{2n^2\varepsilon^2+o(\varepsilon^2)} \\
        & = \lim_{\varepsilon \searrow 0} \frac{4-8(1-\varepsilon)^n+4(1-\varepsilon)^{2n}}{2n^2\varepsilon^2+o(\varepsilon^2)} \\ 
        & = \lim_{\varepsilon \searrow 0} \frac{4-8+8n\varepsilon-8\frac{n(n-1)}{2}\varepsilon^2+o(\varepsilon^2)+4-8n\varepsilon+4\frac{2n(2n-1)}{2}\varepsilon^2+o(\varepsilon^2)}{2n^2\varepsilon^2+o(\varepsilon^2)} \\
        & = \lim_{\varepsilon \searrow 0} \frac{4n^2\varepsilon^2+o(\varepsilon^2)}{2n^2\varepsilon^2+o(\varepsilon^2)} 
         = \lim_{\varepsilon \searrow 0} \frac{4n^2+\frac{o(\varepsilon^2)}{\varepsilon^2}}{2n^2+\frac{o(\varepsilon^2)}{\varepsilon^2}} 
         = \frac{4n^2}{2n^2} = 2.
    \end{align*}
    Note that the same argument applies provided that each space $\X_k$ contains a set of measure $1/2$.
\end{proof}

\end{document}
